\documentclass[11pt]{article}

\usepackage[margin=1in]{geometry}
\usepackage{amsmath}
\usepackage{amssymb}
\usepackage{booktabs}
\usepackage{array}

\title{\textbf{Hoeffding-Based Safety Probability Proposal for Arbitrary Trajectories}}
\author{}
\date{}

\begin{document}
\maketitle

\section*{Objective}

The goal is to validate whether an arbitrary candidate trajectory is safe with a statistical lower bound on the probability of safety.
Instead of proving that a trajectory converges to a nominal reference, the proposal estimates the probability that a candidate trajectory satisfies safety constraints under sampled uncertainty.

\section*{Proposed Validation Artifact}

The proposed validation samples can be stored in:
\[
\texttt{papers/report/arbitrary\_trajectory\_safety\_validation\_placeholder.csv}.
\]
Each row stores one sampled rollout or perturbation test for a candidate trajectory.

\section*{Arbitrary Candidate Trajectory}

Let an arbitrary candidate trajectory be a finite horizon sequence of states and controls.
\[
\tau=
\left\{
(x_0,u_0),(x_1,u_1),\ldots,(x_T,u_T)
\right\}
\]
Plug in the UAV planning interpretation.
\[
x_t=
\begin{bmatrix}
p_{x,t}\\
p_{y,t}\\
p_{z,t}
\end{bmatrix},
\qquad
u_t=
\begin{bmatrix}
v_{x,t}\\
v_{y,t}\\
v_{z,t}
\end{bmatrix}
\]
Final result.
\[
\boxed{\tau \text{ can be any generated, sampled, RRT, LLM, or hand-designed trajectory.}}
\]

\section*{Safety Set}

The trajectory is safe if every state remains inside the workspace and outside all obstacle safety margins.
\[
x_t\in\mathcal X_{\mathrm{safe}},
\qquad
u_t\in\mathcal U,
\qquad
t=0,\ldots,T.
\]
Define the safe set as
\[
\mathcal X_{\mathrm{safe}}
=
\mathcal X
\setminus
\bigcup_{\ell=1}^{N_{\mathrm{obs}}}
\mathcal O_\ell^{r_{\mathrm{safe}}},
\]
where \(\mathcal O_\ell^{r_{\mathrm{safe}}}\) is obstacle \(\ell\) expanded by safety radius \(r_{\mathrm{safe}}\).
Plug in the current planner-style safety margin.
\[
r_{\mathrm{safe}}=0.40\ \mathrm{m}
\]
Final result.
\[
\boxed{
x_t\in\mathcal X_{\mathrm{safe}}
\text{ means the trajectory remains at least }0.40\ \mathrm{m}
\text{ away from obstacles.}
}
\]

\section*{Safety Indicator}

For a sampled rollout \(X_i(\tau)\) of trajectory \(\tau\), define an indicator variable.
\[
I(X_i(\tau)) :=
\begin{cases}
1, & \text{if } x_t\in\mathcal X_{\mathrm{safe}},\ u_t\in\mathcal U,\ \forall t=0,\ldots,T,\\
0, & \text{otherwise.}
\end{cases}
\]
In words, \(I(X_i(\tau))=1\) means the sampled rollout of the arbitrary trajectory is safe.
Plug in one example outcome.
\[
I(X_i(\tau))=1
\quad\Rightarrow\quad
\min_t d(x_t,\mathcal O)\ge0.40\ \mathrm{m}
\]
Final result.
\[
\boxed{I(X_i(\tau))\in[0,1]\text{ is a Bernoulli safety test.}}
\]

\section*{True Safety Probability}

The true probability that trajectory \(\tau\) is safe is
\[
\mu_\tau := \mathbb P[I(X_i(\tau))=1].
\]
This probability is taken over sampled uncertainty such as obstacle noise, tracking error, perception error, wind disturbance, or initial-state perturbation.
Plug in the interpretation.
\[
\mu_\tau
=
\mathbb P[
\tau \text{ remains collision-free and constraint-satisfying}
]
\]
Final result.
\[
\boxed{\mu_\tau\text{ is the safety probability of the arbitrary trajectory.}}
\]

\section*{Empirical Safety Rate}

Given \(p\) independent sampled rollouts of the same arbitrary trajectory, estimate \(\mu_\tau\) by
\[
\tilde\mu_\tau
:=
\frac{1}{p}
\sum_{j=1}^{p}
I(X_j(\tau)).
\]
Plug in an example validation batch with \(p=100\) rollouts and \(96\) safe rollouts.
\[
\tilde\mu_\tau
=
\frac{96}{100}
=0.96
\]
Final result.
\[
\boxed{\tilde\mu_\tau=0.96}
\]

\section*{Hoeffding Inequality}

Let \(I(X_j(\tau))\), \(j=1,\ldots,p\), be \(p\) iid random variables with
\[
0\le I(X_j(\tau))\le1.
\]
Hoeffding's inequality gives
\[
\mathbb P\left[
\left|\tilde\mu_\tau-\mu_\tau\right|
\ge
\epsilon_h
\right]
\le
2\exp(-2p\epsilon_h^2).
\]
Define the confidence level
\[
\delta_h:=2\exp(-2p\epsilon_h^2).
\]
Final result.
\[
\boxed{
\mathbb P\left[
\left|\tilde\mu_\tau-\mu_\tau\right|<\epsilon_h
\right]
\ge
1-\delta_h
}
\]

\section*{Lower Bound on Safety Probability}

Rearrange Hoeffding's inequality to obtain a lower confidence bound.
\[
\mu_\tau
\ge
\tilde\mu_\tau-\epsilon_h
\]
with confidence at least \(1-\delta_h\).
Solve for \(\epsilon_h\) from
\[
\delta_h=2\exp(-2p\epsilon_h^2).
\]
This gives
\[
\epsilon_h
=
\sqrt{
\frac{\ln(2/\delta_h)}{2p}
}.
\]
Final result.
\[
\boxed{
\mu_\tau
\ge
\tilde\mu_\tau
-
\sqrt{
\frac{\ln(2/\delta_h)}{2p}
}
}
\]

\section*{Numerical Example}

Use \(p=100\) sampled rollouts and confidence parameter \(\delta_h=0.05\).
\[
\epsilon_h
=
\sqrt{
\frac{\ln(2/0.05)}{2(100)}
}
\]
\[
\epsilon_h
=
\sqrt{
\frac{\ln(40)}{200}
}
=
\sqrt{
\frac{3.688879}{200}
}
=0.135812
\]
Using \(\tilde\mu_\tau=0.96\), the safety probability lower bound is
\[
\mu_\tau
\ge
0.96-0.135812
=0.824188.
\]
Final result.
\[
\boxed{
\mu_\tau\ge0.824188
\text{ with confidence at least }0.95.
}
\]

\section*{Acceptance Rule}

Choose a required safety probability \(\mu_{\mathrm{crit}}\).
Accept the arbitrary trajectory if
\[
\mu_{\mathrm{crit}}
\le
\tilde\mu_\tau-\epsilon_h.
\]
Plug in \(\mu_{\mathrm{crit}}=0.80\).
\[
0.80
\le
0.96-0.135812
=0.824188
\]
Final result.
\[
\boxed{\tau\text{ is accepted as statistically safe.}}
\]

\section*{Required Number of Rollouts}

If the desired empirical gap is \(\epsilon_h\), solve Hoeffding's inequality for \(p\).
\[
\delta_h=2\exp(-2p\epsilon_h^2)
\]
\[
p
\ge
\frac{\ln(2/\delta_h)}{2\epsilon_h^2}.
\]
Plug in \(\delta_h=0.05\) and \(\epsilon_h=0.05\).
\[
p
\ge
\frac{\ln(40)}{2(0.05)^2}
=
\frac{3.688879}{0.005}
=737.7758
\]
Final result.
\[
\boxed{p\ge738\text{ sampled rollouts are required.}}
\]

\section*{Trajectory-Level Safety Score}

For ranking many candidate trajectories, define the lower confidence bound score
\[
J_{\mathrm{safe}}(\tau)
:=
\tilde\mu_\tau
-
\sqrt{
\frac{\ln(2/\delta_h)}{2p}
}.
\]
The selected trajectory is
\[
\tau^\star
=
\arg\max_{\tau\in\mathcal T_{\mathrm{cand}}}
J_{\mathrm{safe}}(\tau)
\]
subject to
\[
J_{\mathrm{safe}}(\tau)\ge\mu_{\mathrm{crit}}.
\]
Final result.
\[
\boxed{
\text{Choose the trajectory with the largest certified lower bound on safety.}
}
\]

\section*{Closed-Loop Control Integration Ideas}

\subsection*{Idea 1: Safety Filter Before Execution}

Generate candidate trajectories and validate each trajectory with the Hoeffding lower bound.
\[
J_{\mathrm{safe}}(\tau)
\ge
\mu_{\mathrm{crit}}
\]
Only trajectories satisfying the bound are sent to the controller.
Final result.
\[
\boxed{
\text{Unsafe or statistically uncertain trajectories are rejected before flight.}
}
\]

\subsection*{Idea 2: Receding-Horizon Revalidation}

At each control update, validate the next short-horizon trajectory segment.
\[
\tau_{t:t+H}
=
\{x_t,u_t,\ldots,x_{t+H},u_{t+H}\}.
\]
Accept the segment if
\[
J_{\mathrm{safe}}(\tau_{t:t+H})\ge\mu_{\mathrm{crit}}.
\]
Final result.
\[
\boxed{
\text{The safety probability is refreshed online as perception and state estimates change.}
}
\]

\subsection*{Idea 3: Backup Controller Trigger}

If the lower bound drops below the safety threshold, switch to a verified backup controller.
\[
J_{\mathrm{safe}}(\tau)<\mu_{\mathrm{crit}}
\quad\Rightarrow\quad
u=u_{\mathrm{backup}}(x).
\]
The backup controller can hover, brake, land, or track a locally safe waypoint.
Final result.
\[
\boxed{
\text{Closed-loop control uses the probability bound as a runtime safety trigger.}
}
\]

\subsection*{Idea 4: Risk-Aware Cost for MPC}

Include the lower-bound safety score in the MPC objective.
\[
\min_{\tau}
\sum_{t=0}^{H}
\ell(x_t,u_t)
\lambda_{\mathrm{risk}}
\left[
1-J_{\mathrm{safe}}(\tau)
\right]
\]
subject to
\[
J_{\mathrm{safe}}(\tau)\ge\mu_{\mathrm{crit}}.
\]
Final result.
\[
\boxed{
\text{The controller prefers trajectories with higher certified safety probability.}
}
\]

\section*{Summary}

The proposed method turns arbitrary trajectory safety into a Bernoulli validation problem.
For each candidate trajectory, sampled rollouts produce an empirical safety rate \(\tilde\mu_\tau\).
Hoeffding's inequality converts this empirical rate into a lower confidence bound:
\[
\boxed{
\mu_\tau
\ge
\tilde\mu_\tau
-
\sqrt{
\frac{\ln(2/\delta_h)}{2p}
}
}
\]
This bound can be used as a trajectory verifier, a planner ranking score, a receding-horizon safety check, or a trigger for a backup controller.

\end{document}
